\documentclass[runningheads,a4paper]{llncs}

\usepackage{amssymb}
\setcounter{tocdepth}{3}
\usepackage{graphicx}

\usepackage{url}
%\urldef{\mailsa}\path|{alfred.hofmann, ursula.barth, ingrid.haas, frank.holzwarth,|
%\urldef{\mailsb}\path|anna.kramer, leonie.kunz, christine.reiss, nicole.sator,|
%\urldef{\mailsc}\path|erika.siebert-cole, peter.strasser, lncs}@springer.com|    
\newcommand{\keywords}[1]{\par\addvspace\baselineskip
\noindent\keywordname\enspace\ignorespaces#1}

\usepackage[table]{xcolor}
\usepackage{tabularx}
\usepackage{qtree}
\usepackage{color, colortbl}
\usepackage{amsmath}
\usepackage{amssymb}
\usepackage{tikz}
\usepackage{pgfplots}
\usepackage{booktabs}
\usepackage{longtable}
\usepackage{supertabular}
\usepackage{caption}
\usepackage[hidelinks]{hyperref}
\usepackage{IEEEtrantools}
%\usepackage{subfig}
\usepackage{float}
\usepackage{mathtools}
\usepackage{algorithm}
\usepackage{algorithmicx}
\usepackage{algpseudocode}
%\usepackage{commath}
\usepackage{physics}
\usepackage{algpseudocode}
\usepackage{authblk}
\usepackage{hyperref}
\usepackage{eqexpl}
\usepackage{makecell}
\usepackage{subcaption}
%\usepackage{subfigure}
\usepackage{multirow}
\usepackage{soul}
\usepackage{hhline}
\usepackage{stmaryrd}
\usepackage{import}
%\usepackage{filecontents}
\usepackage{catchfile}
%\usepackage{biblatex}
\usepackage[utf8]{inputenc}
\usepackage[english]{babel}
%\usepackage{natbib}
\usepackage{cite}
\usepackage{flexisym}


\usepackage{coffeestains}




\DeclareMathOperator{\capacity}{Cap}
\DeclareMathOperator{\range}{range}
\DeclareMathOperator{\elems}{elems}
%\DeclareMathOperator{\range}{\mathfrak{R}}
\DeclareMathOperator{\HEncoder}{HEncoder}

\renewcommand{\algorithmicrequire}{\textbf{Input:}}
\renewcommand{\algorithmicensure}{\textbf{Output:}}


\pgfplotsset{compat=1.18}

%\bibliographystyle{plain}
%\bibliographystyle{abbrv}
%\bibliographystyle{abbrvnat}
%\bibliographystyle{ieeetr}
%\bibliographystyle{author-year}
%\bibliographystyle{ieeetrnat}
\bibliographystyle{splncs}


%\title{A High Capacity Reversible Data Hiding in Encrypted Images Using Multi-MSB Prediction and Huffman Coding}
%\author{Li-Chiun Lin, Cheng-Ta Huang and Shiuh-Jeng Wang}
%\affil{This is some institute}




\begin{document}
	
	\mainmatter  % start of an individual contribution
	
	% first the title is needed
	\title{
%		Lecture Notes in Computer Science:\\
%		Authors' Instructions for the Preparation\\
%		of Camera-Ready Contributions\\
%		to LNCS/LNAI/LNBI Proceedings
		A High Capacity Reversible Data Hiding in Encrypted Images Using Multi-MSB Prediction and Huffman Coding
	}
	
	% a short form should be given in case it is too long for the running head
	\titlerunning{Reversible Data Hiding Using Multi-MSB Prediction and Huffman Coding}
	
	% the name(s) of the author(s) follow(s) next
	%
	% NB: Chinese authors should write their first names(s) in front of
	% their surnames. This ensures that the names appear correctly in
	% the running heads and the author index.
	%
	\author{
		Li-Chiun Lin
		%
%		\thanks{
%%			Please note that the LNCS Editorial assumes that all authors have used
%%			the western naming convention, with given names preceding surnames. This determines
%%			the structure of the names in the running heads and the author index.
%		}
%
		\and Cheng-Ta Huang 
		\and Shiuh-Jeng Wang
%		\and Frank Holzwarth\and\\
%		Anna Kramer\and Leonie Kunz\and Christine Rei\ss\and\\
%		Nicole Sator\and Erika Siebert-Cole\and Peter Stra\ss er
	}
	%
	\authorrunning{
%		Lecture Notes in Computer Science: Authors' Instructions
	}
	% (feature abused for this document to repeat the title also on left hand pages)
	
	% the affiliations are given next; don't give your e-mail address
	% unless you accept that it will be published
	\institute{
%		Springer-Verlag, Computer Science Editorial,\\
%		Tiergartenstr. 17, 69121 Heidelberg, Germany\\
%		\mailsa\\
%		\mailsb\\
%		\mailsc\\
%		\url{http://www.springer.com/lncs}
	}
	
	%
	% NB: a more complex sample for affiliations and the mapping to the
	% corresponding authors can be found in the file "llncs.dem"
	% (search for the string "\mainmatter" where a contribution starts).
	% "llncs.dem" accompanies the document class "llncs.cls".
	%
	
	\toctitle{Lecture Notes in Computer Science}
	\tocauthor{Authors' Instructions}
	\maketitle
	
	%\coffeestainA{0.9}{0.85}{-25}{5cm}{1.3cm}
	
	\begin{abstract}
		The prosperity of cloud services stimulates the study of reversible data hiding in encrypted images, for which researchers dedicate to find better methods to adopt the new assumption that the ones who provide the cover images, a.k.a.\ the \emph{Conent-owner}, are no longer the same as the ones who hide the secret message into them, a.k.a.\ the \emph{Data-hider}, and the \emph{Data-hider} might not always be allowed to know the actual content of cover image the \emph{Conent-owner} provided, while the one who receives the stego image must be able to restore the exact cover image after extracting the secret.
		A method of such reversible data hiding in encrypted images has been proposed and 
		we propose improvements to enhance its performance and make it more generalized and stable. 
		Experimental results show that our proposal improves its performance even in the worst-case scenario.
	\end{abstract}
	
	%	\tableofcontents
	
	
	\section{Introduction}
\label{sec:intro}

Data hiding is a technique used to prevent the conveyor from knowing the existence of a secret message while transporting it from the sender to the receiver. 
While the cover image is used to embed secret message might inevitably lose some part of its pixels, a subset of the study of data hiding is dedicated to meet the requirement that the cover image must be losslessly recovered after extracting the secret message, namely \emph{reversible data hiding}. 
Furthermore, a subset within it even started to study the benefits of hiding secret messages into encrypted images, i.e., \emph{reversible data hiding in encrypted images}~\cite{Zhang2011, Liu2020, Jain2021, Di2018}.
The schemata of reversible data hiding in encrypted images are usually categorized into the following types. 

In \emph{Reserving Room before Encryption}~\cite{Ma2013, Agrawal2017, Liao2015}, as its name implies, a preprocessing procedure is applied to a cover image before an encryption procedure is applied to it. 
This type of schemata normally produces a \emph{label map} containing critical information for a data hiding procedure. 
This \emph{label map} is then considered as \emph{auxiliary information} and embedded into the front part of the cover image.

In \emph{Vacating Room after Encryption}~\cite{Puech2008, Wu2019, Li2018}, on the contrary, an encryption procedure is applied to a cover image before a post-process procedure can vacate some part of the encrypted cover image for a data hiding procedure.

Though the schemata of \emph{Vacating Room after Encryption} are more concise since the post-process and the data hiding process can be considered as one process, 
those schemata of \emph{Reserving Room before Encryption} provides more capacity because they can exploit the similarity of adjacent pixels while the input of the post process procedure of \emph{Vacating Room after Encryption} is an encrypted cover image, which is equivalent to random bitstreams with bits in the normal distribution.

When discussing reversible data hiding in encrypted images, a use case that comprises of three roles, namely the \emph{Content-owner}, the \emph{Data-hider} and the \emph{Receiver} is commonly adopted~\cite{Qin2015}.
The \emph{Content-owner}, as its name suggested, is responsible for providing cover images, commonly analogies to a cloud service client user.
The responsibility of the \emph{Data-hider} is to embed some secret message into the cover images that the \emph{Content-owner} provided, e.g., a cloud service provider might need to embed authentication data or annotations into the images owned by its clients. 
The \emph{Receiver} could be the same \emph{Data-hider} or any other entities that request the cover images and the embedded information.

Yin et al.\ proposed a reversible data hiding method in encrypted images using adaptive multi-MSB prediction and Huffman coding~\cite{Yin2020}.
Their experimental results demonstrated the superiority of their method over the methods of~\cite{Puteaux2018, Chen2019, Puyang2019, Puteaux2019, Yi2019}.
In this paper
we propose three improvements that increase the embedding capacity and the stability of their method.

The rest of the paper is organized as follows. 
In Sec.\,\ref{sec:ours}, we describe our improvements.
Experimental results to support our proposal are demonstrated in Sec.\,\ref{sec:experimental}.
Section~\ref{sec:conclusion} concludes this paper. 
	%\section{Related Work}
\label{sec:related}


Yin et al.\ proposed a reversible data hiding in encrypted images method base on median edge detector (MED) predictor~\cite{Weinberger2000} and Huffman coding technique~\cite{IRE1952}.
Their method complied with the common use case that consists of the \emph{Content-owner}, the \emph{Data-hider} and the \emph{Receiver}.
The responsibilities of each role are defined as follows.




%\subsection{a little describe here}

\subsection{The Content-owner}

\subsubsection{Preprocess}

\begin{figure}
	\centering
	\begin{tabular}{|c|c|c|c|c|} 
		
		\multicolumn{1}{c}{\phantom{$x_{i-1, j-1}$}} & \multicolumn{1}{c}{\phantom{$x_{i-1, j-1}$}} &
		\multicolumn{1}{c}{\phantom{$x_{i-1, j-1}$}} & \multicolumn{1}{c}{\phantom{$x_{i-1, j-1}$}} & \multicolumn{1}{c}{\phantom{$x_{i-1, j-1}$}}  
		\\
		\hline
		$x_{0, 0}$ & $x_{0, 1}$ & $x_{0, 2}$      & $\dots$    & $x_{0, w}$ \\ \hline
		$x_{1, 0}$ & $\dots$    & $\dots$         & $\dots$    & $\dots$    \\ \hline 
		$x_{2, 0}$ & $\dots$    & 
		\cellcolor{black!10} $x_{i-1, j-1}$ & \cellcolor{black!10}$x_{i - 1, j}$ & $\dots$   \\ \hline 
		$\dots$    & $\dots$    & 
		\cellcolor{black!10}$x_{i, j - 1}$  & \cellcolor{black!20}$x_{i, j}$     & $\dots$    \\ \hline 
		$x_{h, 0}$ & $\dots$    & $\dots$          & $\dots$        & $x_{h, w}$ \\ \hline 
		
		%			\multicolumn{1}{c}{\phantom{$x_{i-1, j-1}$}} & \multicolumn{1}{c}{\phantom{$x_{i-1, j-1}$}} &
		%			\multicolumn{1}{c}{\phantom{$x_{i-1, j-1}$}} & \multicolumn{1}{c}{\phantom{$x_{i-1, j-1}$}} & \multicolumn{1}{c}{\phantom{$x_{i-1, j-1}$}}  
		%			\\
		\multicolumn{5}{c}{{}}\\
		
		\hline
		\rowcolor{black!30}
		\cellcolor{black!30}$N/A$   & $N/A$       & $N/A$       & $\dots$                         & $N/A$       \\ \hline
		\cellcolor{black!30}$N/A$   & $pv_{1, 1}$ & $pv_{1, 2}$ & $\dots$                         & $pv_{1, w}$ \\ \hline
		\cellcolor{black!30}$N/A$   & $pv_{2, 1}$ & $\dots$     & $\dots$                         & $\dots$     \\ \hline 
		\cellcolor{black!30}$\dots$ & $\dots$     & $\dots$     & $\dots$                         & $\dots$     \\ \hline 
		\cellcolor{black!30}$\dots$ & $\dots$     & $\dots$     & \cellcolor{black!20}$pv_{i, j}$ & $\dots$     \\ \hline 
		\cellcolor{black!30}$N/A$   & $pv_{h, 1}$ & $\dots$     & $\dots$                         & $pv_{h, w}$ \\ \hline 
	\end{tabular}
	
	
	
	\caption{The schematic diagram of pixels and its predicted value.}
	\label{fig:xpv}
\end{figure}

For all pixel $x_{i, j} $ in a cover image $ I$, the predicted value $pv$ of median edge detector predictor is defined as Equation~\eqref{eq:pv} and shown in Figure~\ref{fig:xpv}.
The predicted error $pe$ between each pixel and its predicted value is defined as Equation~\eqref{eq:pe}.
Each pixel, except those in the first row and the first column, has a corresponding value $0 \leq t_{i, j} \leq 8$, as defined in Equation~\eqref{eq:t}.
\begin{align}
	\label{eq:pv}
	pv_{i, j} & = 
	\begin{cases}
		\text{undefined}							   & \text{if } i = 0 \vee j = 0 \\
		\max(x_{i - 1, j}, x_{i, j - 1})               & \text{if } x_{i - 1, j - 1} \leq \min(x_{i - 1, j}, x_{i, j - 1}) \\
		\min(x_{i - 1, j}, x_{i, j - 1})               & \text{if } x_{i - 1, j - 1} \geq \max(x_{i - 1, j}, x_{i, j - 1}) \\
		x_{i - 1, j} + x_{i, j - 1} - x_{i - 1, j - 1} & \textrm{otherwise} 
	\end{cases} \\
	%
	\label{eq:pe}
	pe_{i, j} & = 
	\begin{cases}
		\mathrlap{\text{undefined}}
		\phantom{x_{i - 1, j} + x_{i, j - 1} - x_{i - 1, j - 1}}  & \text{if } i = 0 \vee j = 0 \\
		pv_{i, j} \oplus x_{i, j} & \text{if } 1 \leq i < h \wedge 1 \leq j < w \\
	\end{cases} \\
	%
	\label{eq:t}
	t_{i, j} &= 
	\begin{dcases}
		\mathrlap{8} 
		\phantom{x_{i - 1, j} + x_{i, j - 1} - x_{i - 1, j - 1}}&  \text{if } pe_{i, j} = 0       \\
		7 - \lfloor log_2(pe_{i, j}) \rfloor &  \text{if } pe_{i, j} \neq  0
	\end{dcases} 	
\end{align}

Define $T$ as a two dimensional vector containing all $t_{i, j}$ for $ 1 \leq i < h, 1\leq j < w$,
where $h$ and $w$ denote the height and the width of $I$, respectively.

Let $F$ be a vector containing the count of appearance of each possible value in $T$.
\begin{equation}
	\label{eq:f}
	%	F = \{f_v = \#(t = v) \mid t \in T, 0\leq v \leq 8\}
	%	F = \{f_v = \#\{t \in T \mid t = v\}\}, \text{where } 0\leq v \leq 8.
	F = (f_0, \dots, f_8), \qquad \text{where } f_i = \#\{i \in T\}.
\end{equation}

Let $V$ be a monotonically increasing vector consisting of all possible values in $T$ sorted by $F$.
\begin{equation}
	\label{eq:v}
	%	V = [v \in \elems(T) \mid \forall v_i, v_j \in V \colon i \leq j \iff f_{v_i} \leq f_{v_j} ]
	V = (v_0, \dots, v_8), \qquad \text{where } \forall v_i, v_j \in V \colon i \leq j \iff f_{v_i} \leq f_{v_j} 
\end{equation}

Let $cap(v)$ be a function defined for all $v \in V$, representing the capacity each $t_{i, j} = v$ can provide.
\begin{equation}
	\label{eq:cap}
	cap(v) = 
	\begin{cases}
		8     &  \text{if } v = 8       \\
		v + 1 &  \text{if } 0 \leq v \leq 7
	\end{cases} 
\end{equation}

A predefined Huffman tree $H$ is represented as a vector of a bitstream. 
The topology of $H$ is visualized in Figure~\ref{fig:fixtree}.
\begin{equation}
	H = (00, 01, 100, 101, 1100, 1101, 1110, 11110, 11111)
\end{equation}

\begin{figure}
	\scriptsize
	\centering
	\Tree[. [.0 [.0 00 ]
	[.1 01 ]]
	[.1 [.0 [.0 100 ] [.1 101 ] ]
	[.1 [.0 [.0 1100 ] [.1 1101 ] ]
	[.1 [.0 1110 ]
	[.1 [.0 11110 ]
	[.1 11111 ]]]]]]
	\caption{The predefined Huffman tree.}
	\label{fig:fixtree}
\end{figure}

Let $encode$ be a bijection mapping from $V$ onto $H$, and $decode \colon H \to V$ be its reverse. 
\begin{align}
	\label{eq:encode}
	&encode(t) = c
	\\
	\label{eq:decode}
	&decode(c) = t 
	\\
	&\qquad \text{where } \exists i \in [0, 8] : v_i \in V \wedge h_i \in H \wedge v_i = t \wedge h_i = c \notag
\end{align}


Define $payload$ as the differences between the capacity of $v$ and the length of its Huffman code. 
\begin{equation}
	\label{eq:payload}
	payload(v) = cap(v) - \vert encode(v) \vert 
\end{equation}

Let $T^\prime$ be the concatenation of Huffman code of all $t_{i, j} \in T$ and $LEN$ be its expected code length.
\begin{align}
	\label{eq:encodeT1}
	T^\prime &= \bigparallel_{i = 1}^h  { \bigparallel_{j = 1}^w encode(t_{i, j})} \\
	%	T^\prime &= \Arrowvert_{i = 1}^h  { \bigparallel_{j = 1}^w encode(t_{i, j})} \\
	\label{eq:countlen}
	LEN  &= \sum_{v \in V} (f_v \times \lvert encode(v) \rvert)
\end{align}

The Huffman encoding result of \emph{Lena} $512 \times 512$ is given in Table~\ref{tab:msblena} as an example.
The columns are in conformity with their definitions and the column \emph{$\Delta$} represents the result of $f_v \times payload(v)$. 
Each cell in the row \emph{$\sum$} represents the weighted summation of the corresponding column.

\CatchFileDef{\singlelenamsb}{single-res-table-msb-Lena.tex}{}
\begin{table}
	%	\scriptsize
	\centering
	\caption{An example of \emph{Lena} $512 \times 512$.}
	\label{tab:msblena}
	\begin{tabular}{crlcccr}
		\toprule 
		$v$ &   \multicolumn{1}{c}{$f_v$} &     \multicolumn{1}{c}{$encode(v)$}  &  $cap(v)$ &    $\vert encode(v) \vert$    & $payload(v)$  & $\Delta$ \\
		\midrule 
		\singlelenamsb
		\bottomrule
	\end{tabular}
\end{table}


\subsubsection{Encryption and Embedding Auxiliary information}

The cover image $I$ is encrypted with a key stream $K_I$ using Equation~\eqref{eq:enc}.
\begin{equation}
	\label{eq:enc}
	I_e = I \oplus K_I
\end{equation}

Define reference bytes $REF_e$ as the bytes in the first row and the first column of $I_e$.
\begin{equation*}
	%	REF = [x_{i, j} \mid x_{i, j} \in I_e \wedge (i = 0 \vee j = 0)]
	%	REF = (x^e_{0, 0}, \dots, x^e_{0, w - 1}, x^e_{1, 0}, \dots, x^e_{h - 1, 0})
	REF_e = (\underbrace{x^e_{0, 0}, \dots, x^e_{0, w - 1}}_{\text{first row}}, \underbrace{x^e_{1, 0}, \dots, x^e_{h - 1, 0}}_{\text{first column}}), \qquad \text{where } x^e_{i, j} \in I_e
\end{equation*}

Define auxiliary information $S^{aux}$ as illustrated in Figure~\ref{fig:header1}, where $S^{aux\prime}$ represents the first $w + h - 1$ bytes of $S^{aux}$.
The auxiliary information embedding process follows Algorithm~\ref{algo:embed1}. 
Each byte in $REF_e$ is replaced by 8 bits of $S^{aux\prime}$.
The rest of the $S^{aux}$, denoted as $S^{aux\prime\prime}$, is embedded into the main body of $I_e$ using Equation~\eqref{eq:embed1}. 
Each value $t_{i, j} \in T$ determines the capacity of embedding the corresponding $x^e_{i, j} \in I_e$ can provide as defined in Equation~\eqref{eq:cap}. 
\begin{figure}
	%	\scriptsize
	\centering
	
	\begin{tabularx}{\linewidth}{r c cccc}
		\hhline{*{1}{~}*{5}{-}}
		$S^{aux} =$  
		& \multicolumn{1}{|c|}{\cellcolor{black!10}$encode(0) \Vert \cdots \Vert encode(8)$}
		& \multicolumn{1}{|c|}{\cellcolor{black!20}$LEN$}
		& \multicolumn{2}{|c|}{\cellcolor{black!30}$T^\prime$}
		& \multicolumn{1}{|c|}{\cellcolor{black!40}reference}  
		\\ 
		\hhline{*{1}{~}*{5}{-}}
		& \multicolumn{1}{c}{\upbracefill} & \upbracefill& \multicolumn{2}{c}{\upbracefill}&\upbracefill 
		\\
		& \multicolumn{1}{c}{32 bits }
		& 20 bits
		& \multicolumn{2}{c}{$LEN$ bits  }
		& $w + h - 1 $ bytes 
		\\
		
		%		&&&&&\\
		
		\hhline{*{1}{~}*{3}{-}}
		$S^{aux\prime} =$  
		& \multicolumn{1}{|c|}{\cellcolor{black!10}$encode(0) \Vert \cdots \Vert encode(8)$}
		& \multicolumn{1}{|c|}{\cellcolor{black!20}$LEN$}
		& \multicolumn{1}{|c|}{\cellcolor{black!30}partial $ T^\prime$}
		& \phantom{1}  &
		\\ 
		\hhline{*{1}{~}*{3}{-}}
		& \multicolumn{3}{c}{\upbracefill} & &
		\\
		& \multicolumn{3}{c}{$w + h - 1 $ bytes  }
		&&
		\\
		\hhline{*{4}{~}*{2}{-}}
		$S^{aux\prime\prime} =$  
		&&& 
		& \multicolumn{1}{|c|}{\cellcolor{black!30}partial $ T^\prime$}  
		& \multicolumn{1}{|c|}{\cellcolor{black!40}reference}
		\\ 
		\hhline{*{4}{~}*{2}{-}}
		&&&&\multicolumn{2}{c}{\upbracefill}
		\\
		&&&&\multicolumn{2}{c}{$LEN + 32 + 20 $ bits  }
		\\
	\end{tabularx}
	
	\caption{The auxiliary information.}
	\label{fig:header1}
\end{figure}

\begin{align}
	\label{eq:embed1}
	%	Embed(i, j, S^\ast) & = 
	x^{e\prime}_{i, j}  &= 
	\begin{dcases}
		MSB^8(S^\ast)                   & \text{if }  i = 0 \vee j = 0 \vee 7 \leq t_{i, j} \leq  8\\
		%		&\\
		%		\begin{aligned}
		%			&MSB^{(t_{(i, j)} + 1)}(S^\ast) \\
		%			& \quad + LSB^{(7 - t_{i, j})}(x^e_{i, j})
		%		\end{aligned} & \text{if } 0 \leq t_{i, j} < 7 
		\begin{multlined}
			MSB^{(t_{(i, j)} + 1)}(S^\ast) \\
			+ LSB^{(7 - t_{i, j})}(x^e_{i, j})
		\end{multlined} & \text{if } 0 \leq t_{i, j} < 7 
	\end{dcases} 
	%	\intertext{where}  
	\\ 
	& \text{where} \notag \\
	&
	\begin{aligned}
		%		x^e_{i, j} & \in I_e ,  \\
		S^\ast     & = \text{ a subsequence of } S^{aux}  \\
		MSB^{m}(B) & = \sum\nolimits_{i = 8 - m}^7 b_i \times 2^i, m \in [1, 8] \\
		LSB^{n}(B) & = B \bmod 2^n, n \in [1, 8] \\
		%		B &= [b_i \mid b_i \in \{0, 1\}, i \in [0, 7] ].
		B &= (b_0, \dots, b_i), b_i \in \{0, 1\}, i \in [0, 7]
	\end{aligned} \notag
\end{align}

\begin{algorithm}
	\caption{Generate an encrypted image.}
	\label{algo:embed1}
	
	\begin{algorithmic}[1]
		\Require cover image $I$ of size $w \times h$.
		\Require key stream $K_I$.
		\Require pre-defined Huffman coding tree $H$.
		\Ensure label map $T$.
		\Ensure encrypted image $I_e^\prime$.
		\Ensure coordinate $(c_x, c_y)$ to start information hiding.
		
		%		\State $T \gets [t_{i, j} = -1 \mid i \in [1, h), j \in [1, w)]$. 
		%		\State Let $T$ be a 2-dimensional vector of size $w \times h$ with initial values set to $-1$. 
		%		\State $F \gets \{f_t = 0 \mid t \in [0, 8]\}$. 
		\State Initialize $T$.
		\State Initialize $F$.
		
		
		%		\ForAll{$\{x_{i, j} \in I \mid x_{i, j} \notin REF\}$}
		\ForAll{$\{(i, j) \mid 0 < i < h, 0 < j < w\}$}
		\State		Compute $t$ using Equation~\eqref{eq:t}.
		\State		$f_t \gets f_t + 1$.
		\EndFor
		
		%		\State Let $V$ be the elements of $T$ sorted by $f_t$.
		\State Let $V$ be an vector containing the elements of $T$ sorted by $F$.
		\State Define $encode$ as Equation~\eqref{eq:encode}.
		\State $I_e \gets I \oplus K_I$.
		
		%		\State Compute $T^\prime $ and $ LEN$ using Equation~\eqref{eq:encodeT1} and~\eqref{eq:countlen}.
		\State Compute $T^\prime $ using Equation~\eqref{eq:encodeT1}.
		\State Compute $LEN$ using Equation~\eqref{eq:countlen}.
		
		
		\State Allocate $S^{aux}$ as shown in Figure~\ref{fig:header1}.
		\State $I_e^\prime \gets I_e$
		\State Replace $REF_e$ of $I_e^\prime$ with $S^{aux\prime}$.
		
		\ForAll{$\{(i, j) \mid 0 < i < h, 0 < j < w\}$}
		\If			{$S^{aux\prime\prime} = \emptyset$}
		\State			$(c_x, c_y) \gets (i, j)$
		\State			End of algorithm.
		\Else
		\State			Embed $S^{aux\prime\prime}$ into $x^e_{i, j}$ using Equation~\eqref{eq:embed1}.
		\EndIf
		\EndFor
	\end{algorithmic}
\end{algorithm}

\subsection{The Data-hider}

When giving an encrypted image $I_e^\prime$, a coordinate $(c_x, c_y)$ and a label map $T$, 
the \emph{Data-hider} proceeds the information hiding process describing in Algorithm~\ref{algo:hider1}, 
where the secret information $S$ is encrypted by another secret key $K_S$, 
and produces a stego image $I_w$ 
without the knowledge of the cover image $I$.
The encryption process in Equation~\eqref{eq:enc} provides the compartmentalization between the \emph{Content-owner} and the \emph{Data-hider}.

\begin{algorithm}
	\caption{Embed secret information.}
	\label{algo:hider1}
	
	\begin{algorithmic}[1]
		\Require encrypted image $I_e^\prime$ of size $w \times h$.
		\Require coordinate $(c_x, c_y)$ to start information hiding.
		\Require label map $T$.
		\Require secret information $S$.
		\Ensure stego image $I_w$.
		
		%		\ForAll{$\{t_{i, j} \in T \mid (i = c_x \wedge j \in (c_y, w)) \vee (i \in (c_x, h), j \in [1, w))\}$}
		\ForAll{$\{(i, j) \mid (i = c_x \wedge j \in (c_y, w)) \vee (i \in (c_x, h), j \in [1, w))\}$}
		\If			{$S = \emptyset$}
		\State			End of algorithm.
		\Else
		\State			Embed $S$ into $x^e_{i, j}$ using Equation~\eqref{eq:embed1}.
		\EndIf
		\EndFor
	\end{algorithmic}
\end{algorithm}












\subsection{The Receiver}


%\subsubsection{Retrieve Header}
\subsubsection{Extract Secret}
\label{sec:extract}

Let R be a vector containing $(w + h - 1)\times 8$ bits of $REF_w$ of $I_w$, and let $r_i$ represent the $i$-th bit of R.
\begin{equation}
	%	R = [r_i \mid r_i \in \{0, 1\} \wedge 0 \leq i < \vert REF \vert \times 8]
	R = (r_0, \dots, r_{n}), \qquad \text{where } n = ((w + h - 1) \times 8) - 1, r_i \in \{0, 1\}
\end{equation}

The first 32 bits $(r_0   , \dots, r_{31})$ are the Huffman encoded value of 0 to 8 and 
the following 20 bits $(r_{32}, \dots, r_{51})$ contain the length $LEN$ of the entire encoded label map $T^\prime$, as shown in Figure~\ref{fig:header1}.
For the first 32 bits, using the predefined and preshared Huffman tree $H$, they can be divided into 9 unambiguously Huffman codes, each consequently represents the Huffman code of 0 to 8. 
With this information, the receiver can restore $V$ and $decode$ as defined in Equation~\eqref{eq:v} and~\eqref{eq:decode}.
\begin{align}
	\label{eq:r32}
	&(r_0   , \dots, r_{31}) = encode(0) \Vert \cdots \Vert encode(8) \\
	&(r_{32}, \dots, r_{51}) = LEN^{(0)} \Vert \cdots \Vert LEN^{(19)} \\
	\intertext{where $LEN^{(i)}$ denotes the $i$-th bit of $LEN$, such that }  
	\label{eq:len}
	&LEN  = \sum_{i = 0}^{19} r_{(i + 32)} \times 2^i
\end{align}

Let $x^w_{i, j} \in I_w$ be represented as a bit sequence $b_7 \Vert \cdots \Vert b_0$ 
where $b_u$ represents the $u$-th LSB of $x^w_{i, j}$.
The extraction function is defined as Equation~\eqref{eq:extract}.
\begin{align}
	\label{eq:extract}
	Extrac(i, j)  &= 
	\begin{cases}
		b_7 \Vert \cdots \Vert b_0 & \text{if }  7 \leq t_{i, j} \leq 8 \\
		b_7 \Vert \cdots \Vert b_{7 - t_{i, j}} & \text{if }  0 \leq t_{i, j} \leq 6
	\end{cases}
\end{align}

Algorithm~\ref{algo:extractT} describes the detail of how to restore label map $T$ using $decode$ and Equation~\eqref{eq:extract}.
The output $(c_x, c_y)$ are provided to Algorithm~\ref{algo:extractref} to further extract reference bytes $REF_e$ of the encrypted image $I_e$.
Starting from $({c_x}^\prime, {c_y}^\prime)$, the secret bits can be extracted using similar method as described in Algorithm~\ref{algo:extractT} and~\ref{algo:extractref}, and finally be decrypted with $K_S$.



\begin{algorithm}
	\caption{Extract auxiliary information.}
	\label{algo:extractT}
	
	\begin{algorithmic}[1]
		\Require stego image $I_w$ of size $w \times h$.
		\Ensure Huffman coding $decode$.
		\Ensure label map $T$.
		\Ensure coordinate $(c_x, c_y)$ to extract reference bytes $REF_e$.
		\Ensure length $LEN$ of the encoded label map $T^\prime$.
		
		%		\State Let $R$ be a bit secquence containing $REF$ of $I_w$.
		%		\State Let $B$ be a vector containing all bits in $REF_w$ of $I_w$.
		%		\State $R \gets \{x^w_{i, j} \in I_w \mid i = 0 \vee j = 0\}$.
		
		\State Restore 9 Huffman codes from $(r_0 , \dots,  r_{31})$.
		\State Restore $decode$ from the 9 Huffman codes.
		\State Restore $LEN$ from $(r_{32} , \dots,  r_{51})$ using Equation~\eqref{eq:len}.
		
		%		\State $T \gets \{t_{i, j} = -1 \mid i \in [0, h), j \in [0, w)\}$.
		\State Initialize $T$.
		\State $(d_x, d_y) \gets (1, 1)$.	\label{algo:extractT:idxProce}
		\State $buf \gets \emptyset$.
		\ForAll{$\{r_i \in REF_w \mid i > 52\}$}
		\State	$buf \gets buf \Vert r_i$.
		\If		{$buf$ is a valid Huffman code}
		\State 		$t_{(d_x, d_y)} \gets decode(buf)$.		\label{algo:extractT:incdref} \Comment{$t_{(d_x, d_y)} \in T$}
		\State 		$buf \gets \emptyset$.
		%		\State 		$buf \gets buf remove t_{(d_x, d_y)}$.
		\State		Move $(d_x, d_y)$ to next location of $T$.
		%		\State 		$d_y \gets d_y + 1$.
		%		\If			{$d_y = w$}
		%		\State			$d_y \gets 1$.
		%		\State			$d_x \gets d_x + 1$.
		%		\EndIf
		\EndIf
		\EndFor
		
		%		\ForAll{$\{t_{i, j} \mid t_{i, j} \in T, i \in [1, h), j \in [1, w)\}$} \label{algo:extractT:idxExt}
		\ForAll{$\{(i, j) \mid 0 < i < h, 0 < j < w\}$} \label{algo:extractT:idxExt}
		\State 	$buf \gets buf \Vert Extract(i, j)$ by Equation~\eqref{eq:extract}. \Comment{loop invariant}
		\label{algo:extractT:nextij}
		
		\If		{$buf$ is a valid Huffman code}		
		\State 		$t_{(d_x, d_y)} \gets decode(buf)$. \label{algo:extractT:nextd} 
		\State 		$buf \gets \emptyset$.
		\State		Move $(d_x, d_y)$ to next location of $T$.
		\If			{$d_x = h$}
		\State			$(c_x, c_y) \gets (i, j)$
		\State 			End of algorithm.
		\EndIf
		%		\State 		$d_y \gets d_y + 1$.
		%		\If			{$d_y = w$}
		%		\State			$d_y \gets 1$.
		%		\State			$d_x \gets d_x + 1$.
		%		\EndIf
		\EndIf
		\EndFor
	\end{algorithmic}
\end{algorithm}

\begin{algorithm}
	\caption{Extract reference bytes.}
	\label{algo:extractref}
	
	\begin{algorithmic}[1]
		\Require stego image $I_w$ of size $w \times h$.
		\Require label map $T$.
		\Require coordinate $(c_x, c_y)$ to extract reference bytes $REF_e$.
		\Ensure reference bytes $REF_e$ of $I_e$.
		\Ensure coordinate $({c_x}^\prime, {c_y}^\prime)$ to extract secret information.
		
		%		\State Let $REF$ be a byte sequence $[r_i = 0 \mid i \in [0, h + w)]$.
		\State Initialize $REF_e$.
		\State $k \gets 0$.
		%		\ForAll{$\{ t_{i, j} \in T \mid (i = c_x \wedge j \in (c_y, w)) \vee (i \in (c_x, h), j \in [1, w))\}$}
		\ForAll{$\{ (i, j) \mid (i = c_x \wedge j \in (c_y, w)) \vee (i \in (c_x, h), j \in [1, w))\}$}
		\State 		$buf \gets buf \Vert Extract(i, j)$ by Equation~\eqref{eq:extract}.
		\If			{$\left \vert buf \right \vert \geq 8$}
		\State 			$r_k \gets MSB^8(buf)$. \Comment{the $k$-th byte of $REF_e$}
		\State 			Remove 8 bits from $buf$.
		\State 			$k \gets k + 1$.
		\If				{$k = h + w$}
		\State				$({c_x}^\prime, {c_y}^\prime) \gets (i, j)$
		\State 				End of algorithm.
		\EndIf
		\EndIf
		\EndFor
	\end{algorithmic}
\end{algorithm}







\subsubsection{Decrypt and Restore Cover Image}
\label{sec:restore}

After running Algorithm~\ref{algo:extractT} and~\ref{algo:extractref}, the \emph{Receiver} possesses the stego image $I_w$, the label map $T$ and reference bytes $REF_e$ of $I_e$.
Define $I_e^\prime$ as the result of $I_w$ replacing the reference bytes $REF_w$ with $REF_e$.
Each $x_{i, j}^{e\prime} \in I_e^\prime$ is defined as Equation~\eqref{eq:xeprime}.
\begin{align}
	\label{eq:xeprime}
	x_{i, j}^{e\prime} = &
	\begin{cases}
		r_{j}     & \text{ if } i = 0 \\
		r_{w + i}  & \text{ if } j = 0\\
		x_{i, j}     & \text{ otherwise }
	\end{cases} \\
	&\text{where } r_{j}, r_{w + i} \in REF_e \text{ and }  x_{i, j} \in I_w. \notag
\end{align}

Then $I_e^\prime$ is decrypted by key stream $K_I$ the same way as Equation~\eqref{eq:enc} and becomes $I^\prime$.
The cover image $I$ is restored after Equation~\eqref{eq:restore} being applied to all $x^\prime_{i, j} \in I^\prime$.
It must be emphasized that when applying Equation~\eqref{eq:restore} to $x^\prime_{i, j}$, the pixels $x_{i - i, j - 1}, x_{i - 1, j}$ and $x_{i, j -1}$ have already been restored from the reference bytes of $REF_e$ or the former application of Equation~\eqref{eq:restore}, so that the $pv_{i, j}$ are identical as the embedding process.




\begin{equation}
	\label{eq:restore}
	%	Restore(i, j) = 
	x_{i, j} =
	\begin{cases}
		pv_{i, j}                          & \text{if } t_{i, j} = 8 \\
		\\
		\begin{aligned}
			&MSB^{t_{i, j}}(pv_{i, j}) \\
			&\quad + LSB^{(7 - t_{i, j})}(x^\prime_{i, j}) \\
			&\quad + \underbrace{(1 - {x^\prime_{i, j}}^{(7 - t_{i, j})}) \times 2^{(7 - t_{i, j})}}_
			{\text{flip the } (t_{(i, j)}+1) \text{-th bit}} 
		\end{aligned} & \text{if } 0 \leq t_{i, j} \leq 7
	\end{cases}
\end{equation}


An example of preprocessing, encryption, embedding, extraction, decryption, and restoration is illustrated in Figure~\ref{fig:example}.



\begin{figure}
	\scriptsize
	\centering
	
	\begin{tabular}{@{}r@{}r@{}rc|c|c|c|c|c|c|c|c|ll}
		
		
		&& \multicolumn{10}{c}{ } && \\
		&& & \multicolumn{1}{c}{ }& \multicolumn{2}{@{}l}{MSB}&\multicolumn{4}{c}{ }&\multicolumn{2}{r@{}}{LSB} && \\
		&& & \multicolumn{1}{c}{ }& \multicolumn{2}{l}{$b_7$}&\multicolumn{4}{c}{ }&\multicolumn{2}{r}{$b_0$} && \\ 
		\cline{5-12}
		
		&& \multicolumn{10}{c}{ } && \\
		
		
		\cline{5-12}
		&
		
		\multirow{5}{*}{Preprocess $\left \{ \begin{array}{r}
				\\ \\ \\ \\ \\
			\end{array}  \right.$}
		
		& $x_{i, j}$  & = & 1 & 0 & 0 & 1 & 1 & 1 & 0 & 0 && \\ \cline{5-12}
		&& $pv_{i, j}$ & = & 1 & 0 & 0 & 1 & 0 & 1 & 1 & 0 && \\ 
		\hhline{*{4}{~}*{8}{-}}
		&& $pe_{i, j}$ & = 
		& \cellcolor{lightgray}0 &\cellcolor{lightgray}0 &\cellcolor{lightgray}0 &\cellcolor{lightgray}0 
		& \cellcolor{gray}1      & 0 & 1 & 0    && \\
		\cline{5-12}
		&&            &\multicolumn{1}{c}{}        & \multicolumn{4}{@{}c@{}}{\upbracefill}&\multicolumn{4}{c}{ } && \\
		&& $t_{i, j}$ &\multicolumn{1}{@{}c@{}}{=} & \multicolumn{4}{c}{4}&\multicolumn{4}{c}{ }                  && \\
		
		
		&& \multicolumn{10}{c}{ } && \\
		
		
		\cline{5-12}
		&
		\multirow{3}{*}{Encrypt $\left \{ \begin{array}{r}
				\\ \\ \\ 
			\end{array}  \right.$}
		
		& $x_{i, j}$       & = & 1 & 0 & 0 & 1 & 1 & 1 & 0 & 0 && \\ \cline{5-12}
		&& $\oplus \quad K_I$ & = & 1 & 1 & 1 & 1 & 1 & 1 & 1 & 1 && \\ \cline{5-12} 
		&& $x^e_{i, j}$     & = & 0 & 1 & 1 & 0 & 0 & 0 & 1 & 1 && \\ \cline{5-12}
		
		&& \multicolumn{10}{c}{ } && \\
		
		\cline{5-12}
		&
		
		\multirow{5}{*}{Embed $\left \{ \begin{array}{r}
				\\ \\ \\ \\ \\ 
			\end{array}  \right.$}
		
		& $x^e_{i, j}$     & = & 0 & 1 & 1 & 0 & 0 & 0 & 1 & 1 && \\ \cline{5-12}
		&& $t_{i, j}$ & \multicolumn{1}{@{}c@{}}{=} & \multicolumn{4}{c}{4}&\multicolumn{2}{@{}l}{+1}&\multicolumn{2}{c}{ } && \\
		\cline{5-9}
		&& $S^\ast$        & = & 1 & 1 & 0 & 1 & 0 &  \multicolumn{3}{c}{ }     && \\ \cline{5-9}
		
		&& \multicolumn{10}{c}{ } && \\
		
		
		\hhline{*{4}{~}*{8}{-}}
		&& $x^{e\prime}_{i, j}$ & = 
		& \cellcolor{lightgray}1 & \cellcolor{lightgray}1 & \cellcolor{lightgray}0 & \cellcolor{lightgray}1 
		& \cellcolor{gray}0      & 0 & 1 & 1 &    
		
		\multirow{4}{*}{$\left . \begin{array}{r}
				\\ \\ \\ \\
			\end{array}  \right \}$ Extract} 
		
		
		& \\ 
		\hhline{*{4}{~}*{8}{-}}
		&& \multicolumn{10}{c}{ } && \\
		&& $t_{i, j}$ & \multicolumn{1}{@{}c@{}}{=} & \multicolumn{4}{c}{4}&\multicolumn{2}{@{}l}{+1}&\multicolumn{2}{c}{ } && \\
		\cline{5-9}
		&& $S^\ast$        & = & 1 & 1 & 0 & 1 & 0 &  \multicolumn{3}{c}{ }     && \\ \cline{5-9}
		
		&& \multicolumn{10}{c}{ } && \\
		
		
		
		
		\cline{5-12}
		&& $x^{e\prime}_{i, j}$  & = & 1 & 1 & 0 & 1 & 0 & 0 & 1 & 1 &
		
		\multirow{3}{*}{$\left . \begin{array}{r}
				\\ \\ \\
			\end{array}  \right \}$ Decrypt}
		
		& \\ \cline{5-12}
		&& $\oplus \quad K_I$      & = & 1 & 1 & 1 & 1 & 1 & 1 & 1 & 1 && \\ \cline{5-12}
		&& $x^\prime_{i, j}$     & = & 0 & 0 & 1 & 0 & 1 & 1 & 0 & 0 && \\ \cline{5-12}
		
		&& \multicolumn{10}{c}{ } && \\
		
		%		&& & \multicolumn{11}{@{}l}{ \begin{scriptsize}
		%				\makecell[l]{
		%					It must be emphasized that at this moment \\
		%					$x_{i-1, j - 1}, x_{i-1, j}$ and $x_{i, j -1}$ have already been restored.
		%				}
		%		\end{scriptsize} } \\
		
		\hhline{*{4}{~}*{8}{-}}
		&& $pv_{i, j}$    & = 
		& \cellcolor{lightgray}1 & \cellcolor{lightgray}0 & \cellcolor{lightgray}0 & \cellcolor{lightgray}1 
		& \cellcolor{gray}0      & 1 & 1 & 0 &
		
		\multirow{3}{*}{$\left . \begin{array}{r}
				\\ \\ \\ 
			\end{array}  \right \}$ Restore}
		
		& \\
		\hhline{*{4}{~}*{8}{-}}
		&& $t_{i, j}$ & \multicolumn{1}{@{}c@{}}{=} & \multicolumn{4}{c}{4}&\multicolumn{2}{@{}l}{+1}&\multicolumn{2}{c}{ } && \\
		\hhline{*{4}{~}*{8}{-}}
		&& $x_{i, j}$  & = 
		& \cellcolor{lightgray}1 & \cellcolor{lightgray}0 & \cellcolor{lightgray}0 & \cellcolor{lightgray}1 
		& \cellcolor{gray}1 & 1 & 0 & 0 && \\ 
		\hhline{*{4}{~}*{8}{-}}
		
		&&  & \multicolumn{1}{c}{} & \multicolumn{4}{c}{ }&\multicolumn{1}{@{}c@{}}{$\uparrow$}&\multicolumn{3}{c}{ } && \\
		&&  & \multicolumn{1}{c}{} & \multicolumn{4}{c}{ }&\multicolumn{1}{@{}c@{}}{flip}&\multicolumn{3}{c}{ } && \\
		
	\end{tabular}
	
	
	\caption{A visualized example.}
	\label{fig:example}
\end{figure}






	\CatchFileDef{\singlelenashr}{single-res-table-shr-Lena.tex}{}

\section{Proposed Method}
\label{sec:ours}

In this section, we propose a new reversible data hiding in encrypted images method base on~\cite{Yin2020} with three improvements that have a significant increment on embedding rate (ER) while still complying with the typical requirements described in Section~\ref{sec:intro}. 



\subsection{Apply the Huffman tree building algorithm}
\label{sec:improvetree}

We observed that Yin et al.\ introduced their method based on a predefined Huffman tree 
without clarifying the pros and cons behind this design decision. 
We found that this fixed predefined tree topology constrains the true ability of Huffman coding. 
When we apply the standard Huffman coding algorithm
, its full power can be released~\cite{IRE1952, Leeuwen2008a}.

The tree topology of 
\emph{Lena} $512 \times 512$ is shown in Figure~\ref{fig:lenatree} and the statistic data is given in Table~\ref{tab:shrlena}. 
Since in our proposal the tree topology is no longer a preshared knowledge between the \emph{Content-Owner}, the \emph{Data-hider} and the \emph{Receiver}, we have to extend $S^{aux}$ to convey this information. 
 
A Huffman tree of 9 leaves has $2 \times 9 -1 = 17$ nodes, each needs one bit to represent its role, where 1 is defined for an internal node and 0 for a leaf~\cite{Pu2006}. 
Each leaf has a value $0 \leq v \leq 8$, thus four bits are enough to store each leaf value, i.e., $4 \times 9 = 36$ bits. 
Hence we occupy 53 bits in $S^{aux}$ to store the tree topology and leaf values.

\begin{figure}
	\scriptsize
	\Tree[ 
	[.0 [.0 \{00\}\\5 ]
	[.1 [.0 \{010\}\\7 ] [.1 \{011\}\\8 ]]]
	[.1 [.0 [.0 \{100\}\\3 ]
	[.1 [.0 \{1010\}\\2 ]
	[.1 [.0 \{10110\}\\1 ] [.1 \{10111\}\\0 ]]]]
	[.1 [.0 \{110\}\\6 ] [.1 \{111\}\\4 ]]]] 
	\caption{The Huffman coding tree of \emph{Lena} $512 \times 512$.}
	\label{fig:lenatree}
\end{figure}


\begin{table}
	%	\scriptsize
	\centering
	\caption{\emph{Lena} $512 \times 512$ using a different Huffman tree.}
	
	\begin{tabular}{crlcccr}
		\toprule 
		$v$ &   \multicolumn{1}{c}{$f_v$} &     \multicolumn{1}{c}{$encode$}  &  $cap$ &    $\vert encode \vert$    & $payload$  & $\Delta$ \\
		\midrule  
		\singlelenashr
		\bottomrule   
	\end{tabular}
	
	\label{tab:shrlena}
\end{table}

\subsection{Reduce Reference Bytes}
\label{sec:improveref}

In the method of Yin et al.,
the bytes in the first row and the first column, i.e.\ $REF_e$, must be losslessly preserved, and this is the reason why they must concatenate $REF_e$ after $T^\prime$ as part of $S^{aux}$ and embed them before embedding any secret information 
. 

Though this design does not provide any negative effect since $REF_e$ are $neutral$ base on the fact that every byte in $REF_e$ provides 8-bit capacity while also occupying 8 bits to embed itself, with one more byte being defined as $REF_e$, there is one less byte that can be used for information hiding. 

As shown in~\eqref{eq:pe2}, 
for the bytes in the first row, we define their predicted error $pv$ as exclusive-or with the previous byte. 
Similar logic applies to those in the first column. 
\begin{equation}
	\label{eq:pe2}
	pe(i, j) = 
	\begin{cases}
		\text{undefined},             & \text{if } i = 0 \wedge  0 \leq j < |REF_e| \\
		x_{i, j - 1} \oplus x_{i, j}, & \text{if } i = 0 \\
		x_{i - 1, j} \oplus x_{i, j}, & \text{if } j = 0 \\
		pv(i, j) \oplus x_{i, j}, 	  & \text{if } 1 \leq i < h \wedge 1 \leq j < w\,. \\
	\end{cases} 
\end{equation}

Theoretically, the top-left corner pixel $x_{0, 0}$ is the only critical byte that must be losslessly preserved, but for the following reasons, we define $REF_e$ more than this one byte.
Aside from the reason described in Section~\ref{sec:improvetree} that we need 53 bits to store the Huffman tree, 
more importantly, we need some more bytes in $REF_e$ to avoid violating a loop invariant during 
their information extraction process~\cite{Furia2014}.

There are two indexes in 
their information extraction process, 
i.e., the $(d_x, d_y)$ 
and the $(i, j)$ for the loop.
For 
their information extraction process
to work properly, there is a never mentioned loop invariant  
that $(i, j) < (d_x, d_y)$ must be true, i.e., the index $(d_x, d_y)$ must be in a more advanced position in $I_w$ than $(i, j)$ for all time.

While $(i, j)$ is increased with every loop iteration, $(d_x, d_y)$ is only increased when the $buf$ 
contains a proper Huffman code, which is affected by the capacity of each proceeded $t_{i, j}$.
If the distribution of $negative$ $t_{i, j}$ has a tendency to centralized in some part of $T$ and slows down the incremental speed of $(d_x, d_y)$, 
their information extraction process
will encounter exception and fail.

Increase $REF_e$ helps to avoid this problem. 
As can be seen 
in their information extraction process,
after extracting the information of Huffman tree, the rest of the $REF_e$ will be consumed to produce more $t_{(d_x, d_y)}$ and put $(d_x, d_y)$ at a more advanced position, thus reducing the chance of violating the loop invariant.
Though defining more $REF_e$ bytes helps to amend this issue, it also decreases the bytes that can be dedicated to information hiding.

Even though the minimum $REF_e$ requirement differs from image to image, for example, the Airplane $512 \times 512$ needs only 1 extra $REF_e$ beside the basic 53 bits used to convey the coding tree, while Baboon $512 \times 512$ asks for 16 more $REF_e$ bytes to safely go through 
their information extraction process, 
this loop invariant only exists in extracting process and remain unnoticed to the \emph{Content-Owner}.
We have to define a fixed and large enough $REF_e$ and preshare this knowledge with all three participators in the use case. 

According to our experiments, defining $REF_e$ equals to or larger than $7 + 63$ bytes is enough for all the test images we can acquire~\cite{UniversityofSouthernCalifornia2020, Bas2011, bows2007}, but we still retain the possibility that it might need to be extended for some more complex cover images.


















\subsection{Remove the redundant length column}
\label{sec:improvesize}


While we inspecting the method of Yin et al., we found the unreasonable design of constraining the size of $LEN$ to 20 bits, which is only suitable for images of size smaller than or equal to $512 \times 512$. 
When using an image of larger size, e.g., $1024 \times 1024$, 
it is understandable that the label map $T$ is of size $1023 \times 1023$ by definition, which is only slightly less than $2^{20}$, then the total $LEN$ of $T^\prime$ will sure be exceeded the size of 20 bits. 
In fact, according to the experimental data published by~\cite{Yin2020}, the image \emph{Man $1024 \times 1024$} has $LEN = 3,121,508$ bits, which requires at least 22 bits to represent it since $log_2(3,121,508) = 21.574$.
This can be amended by allocating enough space in the auxiliary information by~\eqref{eq:floor}.
\begin{equation}
	\label{eq:floor}
	\vert LEN \vert \; \leq \lfloor log_2((w - 1) \times (h - 1) \times C ) \rfloor + 1 , 
\end{equation}
where the constant $C$ is the Shannon entropy~\cite{Shannon1948} representing the lower bound of the number of bits required to encode each value $t \in T$ 
in the worst-case scenario when $F$ is a uniform distribution.

When we further analyze 
their information extraction process, 
we realized that the $LEN$ is redundant information. 
Since for both the secret extracting process
and the cover image restoring process,  
the label map $T$ is essential information, 
and the $LEN$ is a by-product of the process of re-constructing $T$ from $I_e^\prime$ or $I_w$, 
hence do not need to be explicitly conveyed in $S^{aux}$. 
In our experiments, it is safe to remove the column $LEN$ from $S^{aux}$ 
.

	\CatchFileDef{\deltatablesample}{delta-table-2.tex}{}
\CatchFileDef{\singlemanmsb}{single-res-table-msb-Man.tex}{}
\CatchFileDef{\singlemanshr}{single-res-table-shr-Man.tex}{}


\section{Experimental results}
\label{sec:experimental}

\subsection{Commonly used test images}

We apply the method of ours and Yin et al.\ on common test images collected from the USC-SIPI image dataset~\cite{UniversityofSouthernCalifornia2020}.
The experimental results are shown in Table~\ref{tab:common}. 
Within each column, there are three sub-columns representing the data of our method, the method of Yin et al.\ and the differences, respectively.

\begin{table}
	\scriptsize
	\centering
	\caption{Experimental results of benchmark images (Ours/Yin et~al./$\Delta$)}
	\label{tab:common}
	\begin{tabular}{l@{}c  @{\quad}  r@{/}r@{/}r  @{\quad} r@{/}r@{/}r  @{\quad} r@{/}r@{/}r  @{\quad} r@{/}r@{/}r }
		\toprule
		\multicolumn{1}{c}{Image} & Size & 
		\multicolumn{3}{c}{Hdr(bit)}            & \multicolumn{3}{c}{Ref(byte)}& \multicolumn{3}{c}{Capacity(bit)}& \multicolumn{3}{c}{LEN(bit)}      \\ 
		\midrule
		%		\midrule
		\deltatablesample
		\bottomrule
	\end{tabular}
\end{table}

We noticed that in the experimental result of~\cite{Yin2020}, they use the identical sample test images as we acquired from USC-SIPI, except for 2 images, i.e., \emph{Baboon} and \emph{Lena}. 
To avoid ambiguity, we included both images in Table~\ref{tab:common} and postfix ``-Yin'' to the 2 images used by Yin et al.

In the column \emph{LEN}, our improvement described in Sec.\,\ref{sec:improvetree} surely reduces the overall length of $T'$ since Huffman coding is always optimal~\cite{Long2001}.
%
The reason why our method produces larger $T^\prime$ when applied to \emph{Lake, Peppers, Lax} and \emph{Man} is that we defined a larger $T$ in our method than Yin et al.
%
Take \emph{Man} $1024 \times 1024$ for an example, 
it is understandable that $T$ is of size $1023 \times 1023$ when applying the method of Yin et al., 
while it is of size $1024 \times 1024 - (7 + 63)$ when applying our method, 
where we use 7 bytes to store the 53 bits of \emph{Hdr} and 63 bytes are defined as $REF_e$.
%
Since our method defined a larger $T$, it is inevitable that there will be a longer $T'$ for some images, as we have witnessed in the column \emph{LEN}. 
%
But when we further exam the coding result of \emph{Man} $1024 \times 1024$ in Table~\ref{tab:man}, our method still produces a slightly shorter average code length per value $t$ in $T$, which reflects the fact that our method adopted a better Huffman coding tree.

In the column \emph{Ref}, our improvement described in Sec.\,\ref{sec:improveref} introduced almost a thousand bytes that can be used for information hiding for a $512 \times 512$ image, which turns into a contribution of averagely more than five thousand bits in \emph{Capacity}. 
This improvement is proportional to the \emph{Size} of the cover image.


\begin{table}
	%	\scriptsize
	\centering
	\caption{The encoding result of \emph{Man} $1024 \times 1024$.}
	\label{tab:man}
	

	\subcaption{when applying~\cite{Yin2020}.}
	\label{tab:mana}
	\begin{tabular}{crlcccr}
		\toprule 
		$v$ &   \multicolumn{1}{c}{$f_v$} &     \multicolumn{1}{c}{$encode$}  &  $cap$ &    $\vert encode \vert$    & $payload$  & $\Delta$ \\
		\midrule 
		\singlemanmsb
		\bottomrule
		\multicolumn{7}{l}{Average code length per value: $3121506 \divisionsymbol 1046529 = 2.982722$}
	\end{tabular}
	
	
	
	\subcaption{when applying our method.}
	\label{tab:manb}
	\begin{tabular}{crlcccr}
		\toprule 
		$v$ &   \multicolumn{1}{c}{$f_v$} &     \multicolumn{1}{c}{$encode$}  &  $cap$ &    $\vert encode \vert$    & $payload$  & $\Delta$ \\
		\midrule 
		\singlemanshr
		\bottomrule
		\multicolumn{7}{l}{Average code length per value: $3127243 \divisionsymbol 1048506 = 2.982570$}
	\end{tabular}
\end{table}


	\section{Conclusion}
\label{sec:conclusion}

In this paper, we analyzed an existed reversible data hiding method proposed by Yin et al., and we proposed three improvements, within which, two of them have significant incremental in embedding capacity that has been shown in Sec.\,\ref{sec:experimental}, while the last improvement aimed to make the method more stable and generalize so that it can accept cover images of larger size. 

Since the method within our discussion is a reversible data hiding method, 
which can losslessly restore the cover image after extracting secret message, 
and also because the stego image is indeed an encrypted image, 
we omitted the comparison using PSNR and MSE in Sec.\,\ref{sec:experimental} for their values will always be $+\infty$ and $0$.
We focused our comparison on the embedding capacity, i.e., BPP, and the experimental result of applying both methods shows that our proposal is superior even in the worst-case scenario when applying on rough images. 
	
	
	
	
	
	
	\bibliography{shorter_huffman}
	
	%	\appendix
	%	\section{Huffman Coding Tree}
\label{app:buildtree}

The Huffman coding tree generation process of Lena $512 \times 512$.


\begin{scriptsize}
	\centering
	\begin{longtable}{c}
		\label{app:htree}
		\begin{tabular}{ccccccccc} \toprule
%			\cellcolor{lightgray} 9818 & \cellcolor{lightgray}9742 & 15247 & 33246 & 44509 & 53359 & 41758 & 24353 & 29089  \\ 
			\Tree[.\cellcolor{lightgray}{9818} [.0 ] ] & 
			\Tree[.\cellcolor{lightgray}{9742} [.1 ] ] & 
			\Tree[.15247 [.2 ] ] & 
			\Tree[.33246 [.3 ] ] & 
			\Tree[.44509 [.4 ] ] & 
			\Tree[.53359 [.5 ] ] & 
			\Tree[.41758 [.6 ] ] & 
			\Tree[.24353 [.7 ] ] & 
			\Tree[.29089 [.8 ] ]     
		\end{tabular}
		\\
		\begin{tabular}{cccccccc} \toprule
%			\cellcolor{lightgray} 15247 & 33246 & 44509 & 53359 & 41758 & 24353 & 29089 & \cellcolor{lightgray}19560  \\
			\Tree[.\cellcolor{lightgray}{15247} [.2 ] ] & 
			\Tree[.33246 [.3 ] ] & 
			\Tree[.44509 [.4 ] ] & 
			\Tree[.53359 [.5 ] ] & 
			\Tree[.41758 [.6 ] ] & 
			\Tree[.24353 [.7 ] ] & 
			\Tree[.29089 [.8 ] ] & 
			\Tree[.\cellcolor{lightgray}{19560} [.9742 [.1 ]] [.9818 [.0 ] ] ]  
		\end{tabular}
		\\
		\begin{tabular}{   c  c  c  c   c  c  c } \toprule
%			33246 & 44509 & 53359 & 41758 & \cellcolor{lightgray}24353 & \cellcolor{lightgray}29089 & 34807  \\ 
			\Tree[.33246 [.3 ] ] & 
			\Tree[.44509 [.4 ] ] & 
			\Tree[.53359 [.5 ] ] & 
			\Tree[.41758 [.6 ] ] & 
			\Tree[.\cellcolor{lightgray}24353 [.7 ] ] & 
			\Tree[.\cellcolor{lightgray}29089 [.8 ] ] & 
			\Tree[.34807 [.15247 [.2 ] ] [.19560  [.9742 [.1 ]] [.9818 [.0 ] ] ]]
		\end{tabular}
		\\
		\begin{tabular}{   c  c  c   c  c  c } \toprule
%			\cellcolor{lightgray}33246   & 44509 & 53359 & 41758 & \cellcolor{lightgray}34807& 53442  \\
			\Tree[.\cellcolor{lightgray}33246 [.3 ] ] & 
			\Tree[.44509 [.4 ] ] & 
			\Tree[.53359 [.5 ] ] & 
			\Tree[.41758 [.6 ] ] & 
			\Tree[.\cellcolor{lightgray}34807 [.15247 [.2 ] ] [.19560  [.9742 [.1 ]] [.9818 [.0 ] ] ]] & 
			\Tree[.53442 [.24353 [.7 ] ] [.29089 [.8 ] ] ]   
		\end{tabular}
		\\
		\begin{tabular}{    c  c   c  c  c } \toprule
%			\cellcolor{lightgray}44509 & 53359& \cellcolor{lightgray}41758& 53442& 68053  \\ 
			\Tree[.\cellcolor{lightgray}44509 [.4 ] ] & 
			\Tree[.53359 [.5 ] ] & 
			\Tree[.\cellcolor{lightgray}41758 [.6 ] ] & 
			\Tree[.53442 [.24353 [.7 ] ] [.29089 [.8 ] ] ] & 
			\Tree[.68053 [.33246 [.3 ] ] [.34807 [.15247 [.2 ] ] [.19560  [.9742 [.1 ]] [.9818 [.0 ] ] ]]] 
		\end{tabular}
		\\
		\begin{tabular}{    c   c  c  c } \toprule
			\cellcolor{lightgray}53359& \cellcolor{lightgray}53442& 68053 & 86267  \\ 
			\Tree[ [.5 ] ]
			& \Tree[ [.24353 [.7 ] ] [.29089 [.8 ] ] ]  
			& \Tree[ [.33246 [.3 ] ] [.34807 [.15247 [.2 ] ] [.19560  [.9742 [.1 ]] [.9818 [.0 ] ] ]]] 
			& \Tree[ [.41758 [.6 ] ] [.44509 [.4 ] ]]
		\end{tabular}
		\\
		\begin{tabular}{     c  c  c } \toprule
%			\cellcolor{lightgray}68053 & \cellcolor{lightgray}86267& 106801  \\ 
			\Tree[.\cellcolor{lightgray}68053 [.33246 [.3 ] ] 
				[.34807 [.15247 [.2 ] ] [.19560  [.9742 [.1 ]] [.9818 [.0 ] ] ]]] & 
			\Tree[.\cellcolor{lightgray}86267 [.41758 [.6 ] ] [.44509 [.4 ] ]] & 
			\Tree[.106801 [.53359 [.5 ] ] [.53442 [.24353 [.7 ] ] [.29089 [.8 ] ] ] ]
		\end{tabular}
		\\
		\begin{tabular}{   c  c } \toprule
%			\cellcolor{lightgray}106801& \cellcolor{lightgray}154320  \\ 
			\Tree[.\cellcolor{lightgray}106801 [.53359 [.5 ] ] [.53442 [.24353 [.7 ] ] [.29089 [.8 ] ] ] ] & 
			\Tree[.\cellcolor{lightgray}154320 [.68053 [.33246 [.3 ] ] 
				[.34807 [.15247 [.2 ] ] [.19560  [.9742 [.1 ]] [.9818 [.0 ] ] ]]]  
				[.86267 [.41758 [.6 ] ] [.44509 [.4 ] ]]]
		\end{tabular}
		\\
		\begin{tabular}{   c } \toprule
%			\cellcolor{lightgray}261121  \\ 
			\Tree[.261121 
				[.106801 [.53359 [.5 ] ] [.53442 [.24353 [.7 ] ] [.29089 [.8 ] ] ] ] 
				[.154320 [.68053 [.33246 [.3 ] ] [.34807 [.15247 [.2 ] ] 
					[.19560  [.9742 [.1 ]] [.9818 [.0 ] ] ]]]  
					[.86267 [.41758 [.6 ] ] [.44509 [.4 ] ]]]]
			\\ \midrule
			
			\Tree[ 
			[.0 [.0 \{00\}\\5 ]
			[.1 [.0 \{010\}\\7 ] [.1 \{011\}\\8 ]]]
			[.1 [.0 [.0 \{100\}\\3 ]
			[.1 [.0 \{1010\}\\2 ]
			[.1 [.0 \{10110\}\\1 ] [.1 \{10111\}\\0 ]]]]
			[.1 [.0 \{110\}\\6 ] [.1 \{111\}\\4 ]]]] 
		\end{tabular}
	\end{longtable}
\end{scriptsize}


%\begin{scriptsize}
%	\centering
%	\begin{longtable}{c}
%%		\label{app:htree}
%		\begin{tabular}{  c  c  c  c  c  c  c  c  c } \toprule
%			\cellcolor{lightgray} 9818 & \cellcolor{lightgray}9742 & 15247 & 33246 & 44509 & 53359 & 41758 & 24353 & 29089 
%			\\ 
%			\Tree[ [.0 ] ]  
%			& \Tree[ [.1 ] ]  
%			& \Tree[ [.2 ] ] 
%			& \Tree[ [.3 ] ] 
%			& \Tree[ [.4 ] ] 
%			& \Tree[ [.5 ] ] 
%			& \Tree[ [.6 ] ] 
%			& \Tree[ [.7 ] ] 
%			& \Tree[ [.8 ] ]     
%		\end{tabular}
%		\\
%		\begin{tabular}{   c  c  c  c  c  c  c  c } \toprule
%			\cellcolor{lightgray} 15247 & 33246 & 44509 & 53359 & 41758 & 24353 & 29089 & \cellcolor{lightgray}19560 
%			\\
%			\Tree[ [.2 ] ]
%			& \Tree[ [.3 ] ]
%			& \Tree[ [.4 ] ]
%			& \Tree[ [.5 ] ]
%			& \Tree[ [.6 ] ]
%			& \Tree[ [.7 ] ]
%			& \Tree[ [.8 ] ]  
%			& \Tree[ [.9742 [.1 ]] [.9818 [.0 ] ] ]  
%		\end{tabular}
%		\\
%		\begin{tabular}{   c  c  c  c   c  c  c } \toprule
%			33246 & 44509 & 53359 & 41758 & \cellcolor{lightgray}24353 & \cellcolor{lightgray}29089 & 34807 
%			\\ 
%			\Tree[ [.3 ] ]
%			& \Tree[ [.4 ] ]
%			& \Tree[ [.5 ] ]
%			& \Tree[ [.6 ] ]
%			& \Tree[ [.7 ] ]
%			& \Tree[ [.8 ] ]   
%			& \Tree[ [.15247 [.2 ] ] [.19560  [.9742 [.1 ]] [.9818 [.0 ] ] ]]
%		\end{tabular}
%		\\
%		\begin{tabular}{   c  c  c   c  c  c } \toprule
%			\cellcolor{lightgray}33246   & 44509 & 53359 & 41758 & \cellcolor{lightgray}34807& 53442 
%			\\
%			\Tree[ [.3 ] ]
%			& \Tree[ [.4 ] ]
%			& \Tree[ [.5 ] ]
%			& \Tree[ [.6 ] ]
%			& \Tree[ [.15247 [.2 ] ] [.19560  [.9742 [.1 ]] [.9818 [.0 ] ] ]]
%			& \Tree[ [.24353 [.7 ] ] [.29089 [.8 ] ] ]   
%		\end{tabular}
%		\\
%		\begin{tabular}{    c  c   c  c  c } \toprule
%			\cellcolor{lightgray}44509 & 53359& \cellcolor{lightgray}41758& 53442& 68053 
%			\\ 
%			\Tree[ [.4 ] ]
%			& \Tree[ [.5 ] ]
%			& \Tree[ [.6 ] ]
%			& \Tree[ [.24353 [.7 ] ] [.29089 [.8 ] ] ]  
%			& \Tree[ [.33246 [.3 ] ] [.34807 [.15247 [.2 ] ] [.19560  [.9742 [.1 ]] [.9818 [.0 ] ] ]]] 
%		\end{tabular}
%		\\
%		\begin{tabular}{    c   c  c  c } \toprule
%			\cellcolor{lightgray}53359& \cellcolor{lightgray}53442& 68053 & 86267 
%			\\ 
%			\Tree[ [.5 ] ]
%			& \Tree[ [.24353 [.7 ] ] [.29089 [.8 ] ] ]  
%			& \Tree[ [.33246 [.3 ] ] [.34807 [.15247 [.2 ] ] [.19560  [.9742 [.1 ]] [.9818 [.0 ] ] ]]] 
%			& \Tree[ [.41758 [.6 ] ] [.44509 [.4 ] ]]
%		\end{tabular}
%		\\
%		\begin{tabular}{     c  c  c } \toprule
%			\cellcolor{lightgray}68053 & \cellcolor{lightgray}86267& 106801 
%			\\ 
%			\Tree[ [.33246 [.3 ] ] [.34807 [.15247 [.2 ] ] [.19560  [.9742 [.1 ]] [.9818 [.0 ] ] ]]] 
%			& \Tree[ [.41758 [.6 ] ] [.44509 [.4 ] ]]
%			& \Tree[ [.53359 [.5 ] ] [.53442 [.24353 [.7 ] ] [.29089 [.8 ] ] ] ]
%		\end{tabular}
%		\\
%		\begin{tabular}{   c  c } \toprule
%			\cellcolor{lightgray}106801& \cellcolor{lightgray}154320 
%			\\ 
%			\Tree[ [.53359 [.5 ] ] [.53442 [.24353 [.7 ] ] [.29089 [.8 ] ] ] ]
%			& \Tree[ [.68053 [.33246 [.3 ] ] [.34807 [.15247 [.2 ] ] [.19560  [.9742 [.1 ]] [.9818 [.0 ] ] ]]]  [.86267 [.41758 [.6 ] ] [.44509 [.4 ] ]]]
%		\end{tabular}
%		\\
%		\begin{tabular}{   c } \toprule
%			\cellcolor{lightgray}261121 
%			\\ 
%			\Tree[ [.106801 [.53359 [.5 ] ] [.53442 [.24353 [.7 ] ] [.29089 [.8 ] ] ] ] [.154320 [.68053 [.33246 [.3 ] ] [.34807 [.15247 [.2 ] ] [.19560  [.9742 [.1 ]] [.9818 [.0 ] ] ]]]  [.86267 [.41758 [.6 ] ] [.44509 [.4 ] ]]]]
%			\\ \midrule
%			
%			\Tree[ 
%				[.0 [.0 \{00\}\\5 ]
%				[.1 [.0 \{010\}\\7 ] [.1 \{011\}\\8 ]]]
%				[.1 [.0 [.0 \{100\}\\3 ]
%				[.1 [.0 \{1010\}\\2 ]
%				[.1 [.0 \{10110\}\\1 ] [.1 \{10111\}\\0 ]]]]
%				[.1 [.0 \{110\}\\6 ] [.1 \{111\}\\4 ]]]] 
%		\end{tabular}
%	\end{longtable}
%\end{scriptsize}
	%	\section{Huffman Tree Serialization}

\subsection{Serialization}


The serialization process of the Huffman tree generated from Lena $512 \times 512$.

\begin{scriptsize}
	\centering
	\begin{longtable}{cc}
		\label{app:serial}
		\begin{tabular}{l}
			\Tree[.$\blacktriangleright\circ\phantom{\blacktriangleright}$ [.$\circ$ 5 
			[.$\circ$ 7  8 ]]
			[.$\circ$ [.$\circ$ 3 
			[.$\circ$  2 
			[.$\circ$  1  0 ]]]
			[.$\circ$  6  4 ]]]
			\\
			\underline{1} \underline{ } \underline{ } \underline{ }
			\underline{ } \underline{ } \underline{ } \underline{ }
			\underline{ } \underline{ } \underline{ } \underline{ }
			\underline{ } \underline{ } \underline{ } \underline{ }
			\underline{ }
			\\
			\underline{ } \underline{ } \underline{ }
			\underline{ } \underline{ } \underline{ }
			\underline{ } \underline{ } \underline{ }
		\end{tabular}
		&
		\begin{tabular}{l}
			\Tree[.$\circ$ [.$\blacktriangleright\circ\phantom{\blacktriangleright}$ 5 
			[.$\circ$ 7  8 ]]
			[.$\circ$ [.$\circ$ 3 
			[.$\circ$  2 
			[.$\circ$  1  0 ]]]
			[.$\circ$  6  4 ]]]
			\\
			\underline{1} \underline{1} \underline{ } \underline{ }
			\underline{ } \underline{ } \underline{ } \underline{ }
			\underline{ } \underline{ } \underline{ } \underline{ }
			\underline{ } \underline{ } \underline{ } \underline{ }
			\underline{ }
			\\
			\underline{ } \underline{ } \underline{ }
			\underline{ } \underline{ } \underline{ }
			\underline{ } \underline{ } \underline{ }
		\end{tabular}
		\\
		\begin{tabular}{l}
			\Tree[.$\circ$ [.$\circ$ {$\blacktriangleright 5 \phantom{\blacktriangleright}$} 
			[.$\circ$ 7  8 ]]
			[.$\circ$ [.$\circ$ 3 
			[.$\circ$  2 
			[.$\circ$  1  0 ]]]
			[.$\circ$  6  4 ]]]
			\\
			\underline{1} \underline{1} \underline{0} \underline{ }
			\underline{ } \underline{ } \underline{ } \underline{ }
			\underline{ } \underline{ } \underline{ } \underline{ }
			\underline{ } \underline{ } \underline{ } \underline{ }
			\underline{ }
			\\
			\underline{0101} \underline{ } \underline{ }
			\underline{ } \underline{ } \underline{ }
			\underline{ } \underline{ } \underline{ }
		\end{tabular}
		&
		\begin{tabular}{l}
			\Tree[.$\circ$ [.$\circ$ 5 
			[.$\blacktriangleright\circ\phantom{\blacktriangleright}$ 7  8 ]]
			[.$\circ$ [.$\circ$ 3 
			[.$\circ$  2 
			[.$\circ$  1  0 ]]]
			[.$\circ$  6  4 ]]]
			\\
			\underline{1} \underline{1} \underline{0} \underline{1}
			\underline{ } \underline{ } \underline{ } \underline{ }
			\underline{ } \underline{ } \underline{ } \underline{ }
			\underline{ } \underline{ } \underline{ } \underline{ }
			\underline{ }
			\\
			\underline{0101} \underline{ } \underline{ }
			\underline{ } \underline{ } \underline{ }
			\underline{ } \underline{ } \underline{ }
		\end{tabular}
		\\
		\begin{tabular}{l}
			\Tree[.$\circ$ [.$\circ$ 5 
			[.$\circ$ {$\blacktriangleright 7 \phantom{\blacktriangleright}$} 8 ]]
			[.$\circ$ [.$\circ$ 3 
			[.$\circ$  2 
			[.$\circ$  1  0 ]]]
			[.$\circ$  6  4 ]]]
			\\
			\underline{1} \underline{1} \underline{0} \underline{1}
			\underline{0} \underline{ } \underline{ } \underline{ }
			\underline{ } \underline{ } \underline{ } \underline{ }
			\underline{ } \underline{ } \underline{ } \underline{ }
			\underline{ }
			\\
			\underline{0101} \underline{0111} \underline{ }
			\underline{ } \underline{ } \underline{ }
			\underline{ } \underline{ } \underline{ }
		\end{tabular}
		&
		\begin{tabular}{l}
			\Tree[.$\circ$ [.$\circ$ 5 
			[.$\circ$ 7 {$\blacktriangleright 8 \phantom{\blacktriangleright}$} ]]
			[.$\circ$ [.$\circ$ 3 
			[.$\circ$  2 
			[.$\circ$  1  0 ]]]
			[.$\circ$  6  4 ]]]
			\\
			\underline{1} \underline{1} \underline{0} \underline{1}
			\underline{0} \underline{0} \underline{ } \underline{ }
			\underline{ } \underline{ } \underline{ } \underline{ }
			\underline{ } \underline{ } \underline{ } \underline{ }
			\underline{ }
			\\
			\underline{0101} \underline{0111} \underline{1000}
			\underline{ } \underline{ } \underline{ }
			\underline{ } \underline{ } \underline{ }
		\end{tabular}
		\\
		\begin{tabular}{l}
			\Tree[.$\circ$ [.$\circ$ 5 
			[.$\circ$ 7 8 ]]
			[.$\blacktriangleright\circ\phantom{\blacktriangleright}$ [.$\circ$ 3 
			[.$\circ$  2 
			[.$\circ$  1  0 ]]]
			[.$\circ$  6  4 ]]]
			\\
			\underline{1} \underline{1} \underline{0} \underline{1}
			\underline{0} \underline{0} \underline{1} \underline{ }
			\underline{ } \underline{ } \underline{ } \underline{ }
			\underline{ } \underline{ } \underline{ } \underline{ }
			\underline{ }
			\\
			\underline{0101} \underline{0111} \underline{1000}
			\underline{ } \underline{ } \underline{ }
			\underline{ } \underline{ } \underline{ }
		\end{tabular}
		&
		\begin{tabular}{l}
			\Tree[.$\circ$ [.$\circ$ 5 
			[.$\circ$ 7 8 ]]
			[.$\circ$ [.$\blacktriangleright\circ\phantom{\blacktriangleright}$ 3 
			[.$\circ$  2 
			[.$\circ$  1  0 ]]]
			[.$\circ$  6  4 ]]]
			\\
			\underline{1} \underline{1} \underline{0} \underline{1}
			\underline{0} \underline{0} \underline{1} \underline{1}
			\underline{ } \underline{ } \underline{ } \underline{ }
			\underline{ } \underline{ } \underline{ } \underline{ }
			\underline{ }
			\\
			\underline{0101} \underline{0111} \underline{1000}
			\underline{ } \underline{ } \underline{ }
			\underline{ } \underline{ } \underline{ }
		\end{tabular}
		\\
		\begin{tabular}{l}
			\Tree[.$\circ$ [.$\circ$ 5 
			[.$\circ$ 7 8 ]]
			[.$\circ$ [.$\circ$ {$\blacktriangleright 3 \phantom{\blacktriangleright}$} 
			[.$\circ$  2 
			[.$\circ$  1  0 ]]]
			[.$\circ$  6  4 ]]]
			\\
			\underline{1} \underline{1} \underline{0} \underline{1}
			\underline{0} \underline{0} \underline{1} \underline{1}
			\underline{0} \underline{ } \underline{ } \underline{ }
			\underline{ } \underline{ } \underline{ } \underline{ }
			\underline{ }
			\\
			\underline{0101} \underline{0111} \underline{1000}
			\underline{0011} \underline{ } \underline{ }
			\underline{ } \underline{ } \underline{ }
		\end{tabular}
		&
		\begin{tabular}{l}
			\Tree[.$\circ$ [.$\circ$ 5 
			[.$\circ$ 7 8 ]]
			[.$\circ$ [.$\circ$ 3 
			[.$\blacktriangleright\circ\phantom{\blacktriangleright}$  2 
			[.$\circ$  1  0 ]]]
			[.$\circ$  6  4 ]]]
			\\
			\underline{1} \underline{1} \underline{0} \underline{1}
			\underline{0} \underline{0} \underline{1} \underline{1}
			\underline{0} \underline{1} \underline{ } \underline{ }
			\underline{ } \underline{ } \underline{ } \underline{ }
			\underline{ }
			\\
			\underline{0101} \underline{0111} \underline{1000}
			\underline{0011} \underline{ } \underline{ }
			\underline{ } \underline{ } \underline{ }
		\end{tabular}
		\\
		\begin{tabular}{l}
			\Tree[.$\circ$ [.$\circ$ 5 
			[.$\circ$ 7 8 ]]
			[.$\circ$ [.$\circ$ 3 
			[.$\circ$  {$\blacktriangleright 2 \phantom{\blacktriangleright}$} 
			[.$\circ$  1  0 ]]]
			[.$\circ$  6  4 ]]]
			\\
			\underline{1} \underline{1} \underline{0} \underline{1}
			\underline{0} \underline{0} \underline{1} \underline{1}
			\underline{0} \underline{1} \underline{0} \underline{ }
			\underline{ } \underline{ } \underline{ } \underline{ }
			\underline{ }
			\\
			\underline{0101} \underline{0111} \underline{1000}
			\underline{0011} \underline{0010} \underline{ }
			\underline{ } \underline{ } \underline{ }
		\end{tabular}
		&
		\begin{tabular}{l}
			\Tree[.$\circ$ [.$\circ$ 5 
			[.$\circ$ 7 8 ]]
			[.$\circ$ [.$\circ$ 3 
			[.$\circ$  2 
			[.$\blacktriangleright\circ\phantom{\blacktriangleright}$  1  0 ]]]
			[.$\circ$  6  4 ]]]
			\\
			\underline{1} \underline{1} \underline{0} \underline{1}
			\underline{0} \underline{0} \underline{1} \underline{1}
			\underline{0} \underline{1} \underline{0} \underline{1}
			\underline{ } \underline{ } \underline{ } \underline{ }
			\underline{ }
			\\
			\underline{0101} \underline{0111} \underline{1000}
			\underline{0011} \underline{0010} \underline{ }
			\underline{ } \underline{ } \underline{ }
		\end{tabular}
		\\
		\begin{tabular}{l}
			\Tree[.$\circ$ [.$\circ$ 5 
			[.$\circ$ 7 8 ]]
			[.$\circ$ [.$\circ$ 3 
			[.$\circ$  2
			[.$\circ$  {$\blacktriangleright 1 \phantom{\blacktriangleright}$} 0 ]]]
			[.$\circ$  6  4 ]]]
			\\
			\underline{1} \underline{1} \underline{0} \underline{1}
			\underline{0} \underline{0} \underline{1} \underline{1}
			\underline{0} \underline{1} \underline{0} \underline{1}
			\underline{0} \underline{ } \underline{ } \underline{ }
			\underline{ }
			\\
			\underline{0101} \underline{0111} \underline{1000}
			\underline{0011} \underline{0010} \underline{0001}
			\underline{ } \underline{ } \underline{ }
		\end{tabular}
		&
		\begin{tabular}{l}
			\Tree[.$\circ$ [.$\circ$ 5 
			[.$\circ$ 7 8 ]]
			[.$\circ$ [.$\circ$ 3 
			[.$\circ$  2
			[.$\circ$  1 {$\blacktriangleright 0 \phantom{\blacktriangleright}$} ]]]
			[.$\circ$  6  4 ]]]
			\\
			\underline{1} \underline{1} \underline{0} \underline{1}
			\underline{0} \underline{0} \underline{1} \underline{1}
			\underline{0} \underline{1} \underline{0} \underline{1}
			\underline{0} \underline{0} \underline{ } \underline{ }
			\underline{ }
			\\
			\underline{0101} \underline{0111} \underline{1000}
			\underline{0011} \underline{0010} \underline{0001}
			\underline{0000} \underline{ } \underline{ }
		\end{tabular}
		\\
		\begin{tabular}{l}
			\Tree[.$\circ$ [.$\circ$ 5 
			[.$\circ$ 7 8 ]]
			[.$\circ$ [.$\circ$ 3 
			[.$\circ$  2 
			[.$\circ$  1  0 ]]]
			[.$\blacktriangleright\circ\phantom{\blacktriangleright}$  6  4 ]]]
			\\
			\underline{1} \underline{1} \underline{0} \underline{1}
			\underline{0} \underline{0} \underline{1} \underline{1}
			\underline{0} \underline{1} \underline{0} \underline{1}
			\underline{0} \underline{0} \underline{1} \underline{ }
			\underline{ }
			\\
			\underline{0101} \underline{0111} \underline{1000}
			\underline{0011} \underline{0010} \underline{0001}
			\underline{0000} \underline{ } \underline{ }
		\end{tabular}
		&
		\begin{tabular}{l}
			\Tree[.$\circ$ [.$\circ$ 5 
			[.$\circ$ 7 8 ]]
			[.$\circ$ [.$\circ$ 3 
			[.$\circ$  2
			[.$\circ$  1 0 ]]]
			[.$\circ$  {$\blacktriangleright 6 \phantom{\blacktriangleright}$}  4 ]]]
			\\
			\underline{1} \underline{1} \underline{0} \underline{1}
			\underline{0} \underline{0} \underline{1} \underline{1}
			\underline{0} \underline{1} \underline{0} \underline{1}
			\underline{0} \underline{0} \underline{1} \underline{0}
			\underline{ }
			\\
			\underline{0101} \underline{0111} \underline{1000}
			\underline{0011} \underline{0010} \underline{0001}
			\underline{0000} \underline{0110} \underline{ }
		\end{tabular}
		\\
		\begin{tabular}{l}
			\Tree[.$\circ$ [.$\circ$ 5 
			[.$\circ$ 7 8 ]]
			[.$\circ$ [.$\circ$ 3 
			[.$\circ$  2
			[.$\circ$  1 0 ]]]
			[.$\circ$  6 {$\blacktriangleright 4 \phantom{\blacktriangleright}$}  ]]]
			\\
			\underline{1} \underline{1} \underline{0} \underline{1}
			\underline{0} \underline{0} \underline{1} \underline{1}
			\underline{0} \underline{1} \underline{0} \underline{1}
			\underline{0} \underline{0} \underline{1} \underline{0}
			\underline{0}
			\\
			\underline{0101} \underline{0111} \underline{1000}
			\underline{0011} \underline{0010} \underline{0001}
			\underline{0000} \underline{0110} \underline{0100}
		\end{tabular}
		& \\
	\end{longtable}
\end{scriptsize}
























\subsection{Unserialization}

The unserialization process of the Huffman tree generated from Lena $512 \times 512$.


\begin{scriptsize}
	\centering
	\begin{longtable}{cc}
		\label{app:unserial}
		\begin{tabular}[t]{l}
			\underline{1} {1} {0} {1} {0} {0} {1} {1}
			{0} {1} {0} {1} {0} {0} {1} {0} {0}
			\\
			{0101} {0111} {1000} {0011} {0010} {0001} {0000} {0110} {0100}
			\\
			\Tree[.$\blacktriangleright\circ\phantom{\blacktriangleright}$ \underline{ } \underline{ } ]
		\end{tabular} 
		&
		\begin{tabular}[t]{l}
			{1} \underline{1} {0} {1} {0} {0} {1} {1}
			{0} {1} {0} {1} {0} {0} {1} {0} {0}
			\\
			{0101} {0111} {1000} {0011} {0010} {0001} {0000} {0110} {0100}
			\\
			\Tree[.$\circ$ [.$\blacktriangleright\circ\phantom{\blacktriangleright}$ \underline{ } \underline{ } ] \underline{ } ]
		\end{tabular}
		\\ 
		\begin{tabular}[t]{l}
			{1} {1} \underline{0} {1} {0} {0} {1} {1}
			{0} {1} {0} {1} {0} {0} {1} {0} {0}
			\\
			\underline{0101} {0111} {1000} {0011} {0010} {0001} {0000} {0110} {0100}
			\\
			\Tree[.$\circ$ [.$\circ$ {$\blacktriangleright 5 \phantom{\blacktriangleright}$} \underline{ } ] \underline{ } ]
		\end{tabular}
		&
		\begin{tabular}[t]{l}
			{1} {1} {0} \underline{1} {0} {0} {1} {1}
			{0} {1} {0} {1} {0} {0} {1} {0} {0}
			\\
			\underline{0101} {0111} {1000} {0011} {0010} {0001} {0000} {0110} {0100}
			\\
			\Tree[.$\circ$ [.$\circ$ 5 
			[.$\blacktriangleright\circ\phantom{\blacktriangleright}$ \underline{ } \underline{ } ]] \underline{ } ]
		\end{tabular}
		\\
		\begin{tabular}[t]{l}
			{1} {1} {0} {1} \underline{0} {0} {1} {1}
			{0} {1} {0} {1} {0} {0} {1} {0} {0}
			\\
			{0101}\underline{0111} {1000} {0011} {0010} {0001} {0000} {0110} {0100}
			\\
			\Tree[.$\circ$ [.$\circ$ 5 
			[.$\circ$ {$\blacktriangleright 7 \phantom{\blacktriangleright}$} \underline{ } ]] \underline{ } ]
		\end{tabular}
		&
		\begin{tabular}[t]{l}
			{1} {1} {0} {1} {0} \underline{0} {1} {1}
			{0} {1} {0} {1} {0} {0} {1} {0} {0}
			\\
			{0101} {0111}\underline{1000} {0011} {0010} {0001} {0000} {0110} {0100}
			\\
			\Tree[.$\circ$ [.$\circ$ 5 
			[.$\circ$ 7 {$\blacktriangleright 8 \phantom{\blacktriangleright}$} ]] \underline{ } ]
		\end{tabular}
		\\
		\begin{tabular}[t]{l}
			{1} {1} {0} {1} {0} {0} \underline{1} {1}
			{0} {1} {0} {1} {0} {0} {1} {0} {0}
			\\
			{0101} {0111}\underline{1000} {0011} {0010} {0001} {0000} {0110} {0100}
			\\
			\Tree[.$\circ$ [.$\circ$ 5 
			[.$\circ$ 7 8 ]] 
			[.$\blacktriangleright\circ\phantom{\blacktriangleright}$ \underline{ } \underline{ } ] ] 
		\end{tabular}
		&
		\begin{tabular}[t]{l}
			{1} {1} {0} {1} {0} {0} {1} \underline{1}
			{0} {1} {0} {1} {0} {0} {1} {0} {0}
			\\
			{0101} {0111}\underline{1000} {0011} {0010} {0001} {0000} {0110} {0100}
			\\
			\Tree[.$\circ$ [.$\circ$ 5 
			[.$\circ$ 7 8 ]] 
			[.$\circ$ 
			[.$\blacktriangleright\circ\phantom{\blacktriangleright}$ \underline{ } \underline{ } ] \underline{ } ] ] 
		\end{tabular}
		\\
		\begin{tabular}[t]{l}
			{1} {1} {0} {1} {0} {0} {1} {1}
			\underline{0} {1} {0} {1} {0} {0} {1} {0} {0}
			\\
			{0101} {0111} {1000}\underline{0011} {0010} {0001} {0000} {0110} {0100}
			\\
			\Tree[.$\circ$ [.$\circ$ 5 
			[.$\circ$ 7 8 ]] 
			[.$\circ$ 
			[.$\circ$ {$\blacktriangleright 3 \phantom{\blacktriangleright}$} \underline{ } ] \underline{ } ] ] 
		\end{tabular}
		&
		\begin{tabular}[t]{l}
			{1} {1} {0} {1} {0} {0} {1} {1}
			{0} \underline{1} {0} {1} {0} {0} {1} {0} {0}
			\\
			{0101} {0111} {1000}\underline{0011} {0010} {0001} {0000} {0110} {0100}
			\\
			\Tree[.$\circ$ [.$\circ$ 5 
			[.$\circ$ 7 8 ]] 
			[.$\circ$ 
			[.$\circ$ 3 
			[.$\blacktriangleright\circ\phantom{\blacktriangleright}$ \underline{ } \underline{ } ] ] \underline{ } ] ] 
		\end{tabular}
		\\
		\begin{tabular}[t]{l}
			{1} {1} {0} {1} {0} {0} {1} {1}
			{0} {1} \underline{0} {1} {0} {0} {1} {0} {0}
			\\
			{0101} {0111} {1000} {0011}\underline{0010} {0001} {0000} {0110} {0100}
			\\
			\Tree[.$\circ$ [.$\circ$ 5 
			[.$\circ$ 7 8 ]] 
			[.$\circ$ 
			[.$\circ$ 3 
			[.$\circ$ {$\blacktriangleright 2 \phantom{\blacktriangleright}$} \underline{ } ] ] \underline{ } ] ] 
		\end{tabular}
		&
		\begin{tabular}[t]{l}
			{1} {1} {0} {1} {0} {0} {1} {1}
			{0} {1} {0} \underline{1} {0} {0} {1} {0} {0}
			\\
			{0101} {0111} {1000} {0011}\underline{0010} {0001} {0000} {0110} {0100}
			\\
			\Tree[.$\circ$ [.$\circ$ 5 
			[.$\circ$ 7 8 ]] 
			[.$\circ$ 
			[.$\circ$ 3 
			[.$\circ$ 2 
			[.$\blacktriangleright\circ\phantom{\blacktriangleright}$ \underline{ } \underline{ } ] ] ] \underline{ } ] ] 
		\end{tabular}
		\\
		\begin{tabular}[t]{l}
			{1} {1} {0} {1} {0} {0} {1} {1}
			{0} {1} {0} {1} \underline{0} {0} {1} {0} {0}
			\\
			{0101} {0111} {1000} {0011} {0010}\underline{0001} {0000} {0110} {0100}
			\\
			\Tree[.$\circ$ [.$\circ$ 5 
			[.$\circ$ 7 8 ]] 
			[.$\circ$ 
			[.$\circ$ 3 
			[.$\circ$ 2 
			[.$\circ$ {$\blacktriangleright 1 \phantom{\blacktriangleright}$} \underline{ } ] ] ] \underline{ } ] ] 
		\end{tabular}
		&
		\begin{tabular}[t]{l}
			{1} {1} {0} {1} {0} {0} {1} {1}
			{0} {1} {0} {1} {0} \underline{0} {1} {0} {0}
			\\
			{0101} {0111} {1000} {0011} {0010} {0001}\underline{0000} {0110} {0100}
			\\
			\Tree[.$\circ$ [.$\circ$ 5 
			[.$\circ$ 7 8 ]] 
			[.$\circ$ 
			[.$\circ$ 3 
			[.$\circ$ 2 
			[.$\circ$ 1 {$\blacktriangleright 0 \phantom{\blacktriangleright}$} ] ] ] \underline{ } ] ] 
		\end{tabular}
		\\
		\begin{tabular}[t]{l}
			{1} {1} {0} {1} {0} {0} {1} {1}
			{0} {1} {0} {1} {0} {0} \underline{1} {0} {0}
			\\
			{0101} {0111} {1000} {0011} {0010} {0001}\underline{0000} {0110} {0100}
			\\
			\Tree[.$\circ$ [.$\circ$ 5 
			[.$\circ$ 7 8 ]] 
			[.$\circ$ 
			[.$\circ$ 3 
			[.$\circ$ 2 
			[.$\circ$ 1 0 ] ] ] [.$\blacktriangleright\circ\phantom{\blacktriangleright}$ \underline{ } \underline{ } ] ] ] 
		\end{tabular}
		&
		\begin{tabular}[t]{l}
			{1} {1} {0} {1} {0} {0} {1} {1}
			{0} {1} {0} {1} {0} {0} {1} \underline{0} {0}
			\\
			{0101} {0111} {1000} {0011} {0010} {0001} {0000}\underline{0110} {0100}
			\\
			\Tree[.$\circ$ [.$\circ$ 5 
			[.$\circ$ 7 8 ]] 
			[.$\circ$ 
			[.$\circ$ 3 
			[.$\circ$ 2 
			[.$\circ$ 1 0 ] ] ] 
			[.$\circ$ {$\blacktriangleright 6 \phantom{\blacktriangleright}$} \underline{ } ] ] ] 
		\end{tabular}
		\\
		\begin{tabular}[t]{l}
			{1} {1} {0} {1} {0} {0} {1} {1}
			{0} {1} {0} {1} {0} {0} {1} {0} \underline{0}
			\\
			{0101} {0111} {1000} {0011} {0010} {0001} {0000} {0110}\underline{0100}
			\\
			\Tree[.$\circ$ [.$\circ$ 5 
			[.$\circ$ 7 8 ]] 
			[.$\circ$ 
			[.$\circ$ 3 
			[.$\circ$ 2 
			[.$\circ$ 1 0 ] ] ] 
			[.$\circ$ 6 {$\blacktriangleright 4 \phantom{\blacktriangleright}$} ] ] ] 
		\end{tabular}
		& \\
	\end{longtable}
\end{scriptsize}

	
	
	
\end{document}